\subsection{Motivation: Why Innovation Over Gradient Difference?}

Consider optimizing $f(\x)$ with stochastic gradients $\g_t = \nabla f(\x_t) + \xi_t$ where $\E[\xi_t] = 0$, $\Var[\xi_t] = \sigma^2$.

\textbf{Adan's gradient difference:}
\begin{equation}
d_t = \g_t - \g_{t-1}
\end{equation}
This compares \emph{two noisy samples}. Under weak correlation:
\begin{equation}
\Var[d_t] = \Var[\g_t] + \Var[\g_{t-1}] - 2\Cov[\g_t, \g_{t-1}] \approx 2\sigma^2
\end{equation}

\textbf{IRIS's innovation:}
\begin{equation}
\Innovation_{t|t-1} = \g_t - \ghat_{t-1}
\end{equation}
This compares \emph{noisy sample to smooth estimate} $\ghat_{t-1} = \beta_1 \ghat_{t-2} + (1-\beta_1)\g_{t-1}$:
\begin{equation}
\Var[\Innovation_{t|t-1}] = \Var[\g_t] + \underbrace{\Var[\ghat_{t-1}]}_{\approx 0} \approx \sigma^2
\end{equation}

\begin{proposition}[Variance Reduction]
\label{prop:variance}
For stochastic gradients with variance $\sigma^2$ and weak temporal correlation, gradient difference has approximately 2$\times$ the variance of innovation:
\begin{equation}
\frac{\Var[d_t]}{\Var[\Innovation_{t|t-1}]} \approx 2
\end{equation}
\end{proposition}

\textbf{Implication:} Innovation's 50\% lower variance enables stable convergence even with high variance smoothing ($\beta_3 \geq 0.999$), allowing 10-1000$\times$ higher learning rates.

\subsection{The IRIS Algorithm}

IRIS maintains three bias-corrected estimates:
\begin{itemize}
    \item $\ghat_t$: Gradient estimate (momentum)
    \item $\epshat_t$: Innovation residual (systematic error accumulation)
    \item $\vhat_t$: Variance estimate (uncertainty of corrected gradient)
\end{itemize}

\begin{algorithm}
\caption{IRIS: Innovation Residual Iterative Stabilization}
\begin{algorithmic}
\REQUIRE Learning rate $\eta$, weight decay $\lambda$, numerical stability $\varepsilon$
\REQUIRE EMA coefficients $(\beta_1, \beta_2, \beta_3)$, optional split-correction $\gamma$
\STATE Initialize $\widehat{\mathbf{g}}_0 \gets \mathbf{0}, \; \widehat{\bm{\varepsilon}}_0 \gets \mathbf{0}, \; \widehat{\mathbf{v}}_0 \gets \mathbf{0}$
\STATE Initialize bias correction $\psi_{1,0} \gets 0, \; \psi_{2,0} \gets 0, \; \psi_{3,0} \gets 0$
\FOR{$t = 1, 2, 3, \ldots$}
    \STATE $\psi_{i,t} \gets \beta_i \cdot \psi_{i,t-1} + 1, \quad \forall i \in \{1,2,3\}$ \hfill // Bias correction accumulators
    \IF{split-correction $\gamma$}
    \STATE $\varphi_{t} \gets \gamma \cdot \varphi_{t-1} + 1, \quad \forall i \in \{1,2,3\}$ \hfill // Bias correction accumulators
    \STATE $\textcolor{blue}{\mathbf{I}_{t|t-1}} \gets \mathbf{g}_t - \textcolor{blue}{\widehat{\mathbf{g}}_{t-1}}$ \hfill // Innovation (prediction error)
    \STATE $\textcolor{blue}{\widehat{\mathbf{g}}_t} \gets \textcolor{blue}{\widehat{\mathbf{g}}_{t-1}} + \psi_{1,t}^{-1} \textcolor{blue}{\mathbf{I}_{t|t-1}}$ \hfill // Update gradient estimate
    \STATE $\epshat_t\gets(1-\psi_{2,t}^{-1})\epshat_{t-1}+\psi_{2,t}^{-1} \textcolor{blue}{\mathbf{I}_{t|t-1}}$
    \STATE $\textcolor{blue}{\widehat{\mathbf{v}}_t} \gets (1-\psi_{3,t}^{-1})\textcolor{blue}{\widehat{\mathbf{v}}_{t-1}} + \psi_{3,t}^{-1} {\mathbf{g}_{t}^2}$ \hfill // Update gradient estimate
    \STATE $\nhat_t\gets(1-\varphi_{t}^{-1})\nhat_{t-1}+\varphi_{t}^{-1} \textcolor{blue}{\mathbf{I}_{t|t-1}^2}$
    \STATE $\mathbf{x}_t \gets (1 - \eta\lambda) \mathbf{x}_{t-1} - \eta \left(\dfrac{\textcolor{blue}{\widehat{\mathbf{g}}_t} }{\sqrt{\textcolor{blue}{\widehat{\mathbf{v}}_t}} + \varepsilon}+\beta_2\dfrac{{\widehat{\bm{\varepsilon}}_t}}{\sqrt{\widehat{\mathbf{n}}_t} + \varepsilon}\right)$ \hfill // Parameter update
    
    \ELSE
    \STATE $\textcolor{blue}{\mathbf{I}_{t|t-1}} \gets \mathbf{g}_t - \textcolor{blue}{\widehat{\mathbf{g}}_{t-1}}$ \hfill // Innovation (prediction error)
    \STATE $\textcolor{OliveGreen}{\mathbf{E}_{t|t-1}} \gets \textcolor{blue}{\mathbf{I}_{t|t-1}} - \textcolor{OliveGreen}{\widehat{\bm{\varepsilon}}_{t-1}}$ \hfill // Error-of-innovation
    \STATE ${\bm{\Sigma}_{t|t-1}} \gets (\mathbf{g}_t + \beta_2 \textcolor{blue}{\mathbf{I}_{t|t-1}})^2 - \widehat{\mathbf{v}}_{t-1}$ \hfill // Variance innovation
    \STATE $\textcolor{blue}{\widehat{\mathbf{g}}_t} \gets \textcolor{blue}{\widehat{\mathbf{g}}_{t-1}} + \psi_{1,t}^{-1} \textcolor{blue}{\mathbf{I}_{t|t-1}}$ \hfill // Update gradient estimate
    \STATE $\textcolor{OliveGreen}{\widehat{\bm{\varepsilon}}_t} \gets \textcolor{OliveGreen}{\widehat{\bm{\varepsilon}}_{t-1}} + \psi_{2,t}^{-1} \textcolor{OliveGreen}{\mathbf{E}_{t|t-1}}$ \hfill // Update innovation residual
    \STATE $\widehat{\mathbf{v}}_t \gets \widehat{\mathbf{v}}_{t-1} + \psi_{3,t}^{-1} \bm{\Sigma}_{t|t-1}$ \hfill // Update variance
    \STATE $\mathbf{x}_t \gets (1 - \eta\lambda) \mathbf{x}_{t-1} - \eta \dfrac{\textcolor{blue}{\widehat{\mathbf{g}}_t} + \beta_2 \textcolor{OliveGreen}{\widehat{\bm{\varepsilon}}_t}}{\sqrt{\widehat{\mathbf{v}}_t} + \varepsilon}$ \hfill // Parameter update
    \ENDIF
\ENDFOR
\STATE \textbf{return} $\mathbf{x}_t$
\end{algorithmic}
\end{algorithm}

\subsection{Key Design Choices}

\subsubsection{Variance Matching (Line 9)}

The variance term $(\g_t + \beta_2 \Innovation_{t|t-1})^2$ estimates uncertainty of the \emph{corrected gradient}---exactly what appears in the numerator. This ensures the preconditioner scale matches the update scale.

\textbf{Contrast with AdaBelief:} AdaBelief uses $(\g_t - \m_t)^2$ in the \emph{denominator}, causing:
\begin{itemize}
    \item Division by small numbers when $\g_t \approx \m_t$ (instability)
    \item Inverse scaling: small prediction error $\implies$ large step (counterintuitive)
    \item No variance matching: denominator doesn't match numerator structure
\end{itemize}

IRIS uses innovation in the \emph{numerator}: correction modifies \emph{direction}, not scale. Variance matching in denominator provides stable preconditioning.

\subsubsection{Recursive Bias Correction (Line 4)}

Standard Adam uses $\hat{m}_t = m_t / (1 - \beta_1^t)$, requiring:
\begin{itemize}
    \item Integer $t$ from GPU (GPU-CPU synchronization)
    \item Exponentiation $\beta_1^t$ (numerical instability at large $t$)
\end{itemize}

IRIS uses recursive accumulation:
\begin{equation}
\psi_{i,t} = \beta_i \psi_{i,t-1} + 1 \implies \psi_{i,t}^{-1} = \frac{\beta_i^{-1}}{\psi_{i,t-1}^{-1} + \beta_i^{-1}}
\end{equation}

This is equivalent to Kalman gain with static measurement noise $R = \beta_i^{-1}$:
\begin{equation}
K_t = \frac{P_{t-1}}{P_{t-1} + R} \quad \text{(Kalman filter)}
\end{equation}
where $P_t \leftrightarrow \psi_t^{-1}$ (accumulated uncertainty), $R \leftrightarrow \beta^{-1}$ (measurement noise).

\textbf{Benefits:} No GPU-CPU sync, no exponentiation, exact bias correction from step 1 $\implies$ \textbf{no warmup required}.

\subsubsection{Timescale Separation (\texorpdfstring{$\beta_2 < \beta_1$}{beta2 < beta1})}

We empirically find optimal $\beta_2 = \beta_1 - \delta$ where $\delta \in [0.15, 0.25]$. This creates timescale separation:
\begin{align}
\tau_{\text{gradient}} &= \frac{1}{1-\beta_1} \approx 50 \text{ steps (slow: tracks true gradient)} \\
\tau_{\text{error}} &= \frac{1}{1-\beta_2} \approx 12 \text{ steps (fast: reacts to errors)}
\end{align}

\textbf{Intuition:} Error correction must adapt \emph{faster} than the error itself. If $\beta_2 \approx \beta_1$, both move at same timescale $\implies$ always chasing, never catching.

This is hierarchical adaptive control: fast inner loop stabilizes, slow outer loop tracks.

\subsection{Comparison with Adan}

\begin{table}[h]
\centering
\caption{IRIS vs Adan: Algorithmic Comparison}
\label{tab:iris_vs_adan}
\small
\begin{tabular}{@{}lll@{}}
\toprule
\textbf{Property} & \textbf{Adan} & \textbf{IRIS} \\
\midrule
Core quantity & $d_t = \g_t - \g_{t-1}$ & $\Innovation_{t|t-1} = \g_t - \ghat_{t-1}$ \\
Interpretation & Gradient velocity & Prediction error \\
Variance & $2\sigma^2$ (two samples) & $\sigma^2$ (vs smooth estimate) \\
Near-minima behavior & $d_t \to 0$ (correction vanishes) & $\Innovation_{t|t-1} < 0$ (active damping) \\
Stored gradients & Requires $\g_{t-1}$ & No gradient storage \\
Memory buffers & 4 ($\m$, $\mathbf{v}$, $\mathbf{d}$, $\g_{t-1}$) & 3 ($\ghat$, $\vhat$, $\epshat$) \\
Bias correction & Standard ($\beta^t$) & Kalman gain (recursive) \\
Stability & Convergence-divergence cycles & Unconditional stability \\
\bottomrule
\end{tabular}
\end{table}

\textbf{Key difference:} Adan's gradient difference becomes \emph{less informative} as training progresses (approaches zero near minima). IRIS's innovation remains \emph{informative throughout}---near minima, lagging momentum creates large negative innovation that automatically damps updates.
